\section{Historische Einordnung}

\subsection*{Leitfragen zur Recherche}
\begin{itemize}
    \item \textbf{Entwicklung:} Wie entwickelte sich die Anamorphose aus der Zentralperspektive der Renaissance?
    \item \textbf{Protagonisten:} Welche Rolle spielten Leonardo da Vinci (frühe Skizzen) und Jean-François Niceron (Theorie) bei der Entwicklung?
    \item \textbf{Motivation:} Warum nutzten Künstler Anamorphosen? (Verschlüsselung geheimer Botschaften, Memento Mori, wissenschaftliche Spielerei).
    \item \textbf{Abgrenzung:} Was ist der historische Unterschied zwischen optischen Anamorphosen (Blickwinkel) und katoptrischen Anamorphosen (Zylinderspiegel)?
\end{itemize}

% TODO: Hier Fließtext schreiben, der die Fragen beantwortet.
Um uns dem Gegenstand dieser Untersuchung zu nähern, lade ich Sie zunächst zu einem kurzen Gedankenexperiment ein.

Stellen Sie sich als Wanderer auf einem Berg vor, welcher sich während seiner Reise verirrt hat. Die Kiefern vermischen ihren schweren, harzigen Duft mit der Bergluft, was Sie spürbar beruhigt. Ihr Blick schweift vom kargen Berg hinab in das satte Grün des Waldes, welcher unten in einem kesselförmigen Tal liegt. Als Nächstes folgen Sie einem schmalen Steig. Ein dichter, erdiger Duft von feuchtem Moos steht in der Luft. Während Sie dem Pfad weiter folgen, vernehmen Sie aus der Ferne ein klares und helles Klingen: Es ist der Waldbach, dessen kaltes Wasser leise über die kalten Steinplatten gleitet, aber dazwischen einen hellen Ton aufklingen lässt, wenn das Wasser über einen kleinen Kieselstein springt. Der Ton wird klarer, je näher Sie kommen, bis Sie schließlich abseits des Weges am Flussufer stehen.

Abseits des Weges stellen sie fest, dass der Fluss nicht mehr als Zwei Schritte breit ist. Als Sie die Breite des Flusses betrachten, sehen Sie das in der Sonne glänzende Moos, das Sie schon auf dem Hinweg wahrgenommen haben.

Aufgrund Ihres ursprünglichen Plans, den Berg hinabzuwandern, und Ihrer Vermutung, dass das in rasanter Geschwindigkeit fließende Wasser sicher nach unten führt, beschließen Sie, dem Fluss zu folgen.

Über Stock und Stein folgen Sie dem Fluss und stehen schließlich im Tal vor einem Weiher, in den der Fluss über mehrere Adern verteilt mündet. Das Moos und die Steine sind zwischen den Flussadern verteilt. Sie schauen in den Weiher und bemerken, dass dieser besonders klar ist. Durch die straff gespannte Wasseroberfläche des ruhigen Weihers können sie mühelos den Grund erkennen. 

Es macht den Anschein, als würde sich die fast unsichtbare Wasseroberfläche wie ein Kleid an den Grund anschmiegen. Durch feine Zwischenräume verschafft sich Ihr Blick Zugang zu den schimmernden Sedimenten am Grund des Weihers. Sie schauen hindurch und vollziehen damit unbewusst genau das, was der lateinische Begriff perspicere beschreibt.

Während Sie die schimmernden Steine am Grund des Weihers beobachten, bemerken Sie über die Reflexionen auf dem Wasser, dass die Sonne hinter dem Berg untergeht. Sie beobachten, dass die Form der Sonne eher gestaucht ist und einer Ellipse gleicht anstatt einem Kreis. Weil Sie sich unsicher fühlen, drehen Sie sich abwechselnd von der Wasseroberfläche zur Sonne und wieder zurück. Um das Ungleichgewicht Ihrer gewohnten Wahrnehmung auszugleichen, muss Ihr Geist das Bild umformen - ein "Verrücken" ist nötig, was wir als Anamorphose bezeichnen. Das Wort Anamorphose entspringt dem griechischen Wort anamorphosis und bedeutet die Umformung.